% FILE: 01_research_question.tex
% Project: research-prompt-manufactory
% Thesis Role: prompt-manufacturing-and-subagent-command-surface

\section{Research Question}
\label{sec:research-prompt-manufactory-question}

\subsection{Problem Statement}
\label{sec:research-prompt-manufactory-question-problem}

Multi-agent research programs depend on subagent instructions to direct bounded inspection
stations, enforce surface constraints, and authorize artifact emission. In the conventional
regime, these instructions are hand-written natural-language text: authored by a human,
stored as opaque files, and possessing no traceable derivation from the domain model they
purport to implement. A hand-written prompt cannot be audited against the law it claims to
enforce. It cannot carry a receipt. It cannot be proved to be consistent with the ontology
that governs the research program. When a subagent deviates from its chartered surface, no
manufacturing defect is detectable because the prompt was never manufactured --- it was
merely written.

The Prompt Manufactory addresses this defect directly. The research problem is: \textit{Can
subagent command surfaces be manufactured from graph-encoded law such that every emitted
instruction carries provenance, a type classification, a receipt, and a derivation triple
traceable to a specific ontology instance?} If so, the prompt becomes a first-class
manufactured artifact subject to the same proof requirements as any other artifact in a
receipted manufacturing pipeline.

\subsection{Old-Regime Assumption Challenged}
\label{sec:research-prompt-manufactory-question-old-regime}

The old-regime assumption is that prompts are speech acts: the author's intent expressed
in natural language, communicating desired behavior to an agent. This assumption is deeply
embedded in current practice. Prompt engineering is taught as a craft of phrasing and
iteration. Prompt libraries are version-controlled like documentation. There is no concept
of prompt provenance, prompt type classification, or prompt receipt in the conventional
model.

The Prompt Manufactory challenges this assumption on three grounds. First, if the research
program is governed by an ontology encoding workflow law, subagent roles, hook policies,
and checkpoint verdicts, then a hand-written prompt that claims to implement that law has
no verifiable relationship to it. The prompt may contradict the law without detection.
Second, the forbidden-collapse-law ontology (\texttt{forbidden-collapse-law.ttl}) defines
three absolute bans --- \texttt{.ggen} source files, DTO/JSON flattening, and forced ALIVE
verdicts --- which are graph-expressible constraints. A manufactured warrant can carry these
constraints as typed properties; a hand-written prompt cannot. Third, the Van der Aalst
Constitution (governing all research in this foundry) requires that process claims be
provable from event evidence. A prompt that cannot be traced to its source law produces no
event evidence of lawful derivation.

\subsection{Hypothesis}
\label{sec:research-prompt-manufactory-question-hypothesis}

The hypothesis of the Prompt Manufactory research program is: \textit{A SPARQL-plus-Tera
manufacturing pipeline operating over a typed RDF ontology graph can manufacture all
subagent command surfaces required by a multi-agent research workflow, with zero
hand-written prompt text, such that every emitted warrant carries a
\texttt{pm:derivedFrom} triple and a receipt entry in the manufacturing ledger.}

The secondary hypothesis is that this manufacturing discipline imposes useful
constraints: it prevents forbidden collapses (DTO flattening, forced ALIVE, illegal
surface access) by making constraint violations detectable as ontology violations rather
than runtime defects. A subagent instructed by a manufactured warrant that encodes its
owned and forbidden surfaces cannot plausibly claim ignorance of those surfaces --- the
instruction is typed, receipted, and derived from the same law that governs the audit.

\subsection{Success Criteria}
\label{sec:research-prompt-manufactory-question-success}

The Prompt Manufactory defines success through the 11-gate audit declared in
\texttt{GGEN\_PROMPT\_MANUFACTORY\_PARTIAL\_001.md}. An ALIVE verdict requires all 11 gates
to pass. As of the issued checkpoint (2026-06-01), 9 gates pass:

\begin{enumerate}
  \item Ontology files present and parsing without error (8 \texttt{.ttl} files).
  \item Core SPARQL queries execute and return valid results.
  \item \texttt{ggen.toml} syntax valid; proof specimen template renders correctly.
  \item Seven seed programs encoded in \texttt{research-program-law.ttl}.
  \item End-to-end warrant path proven for at least one program instance.
  \item Zero hand-written program prompts (all warrants carry \texttt{pm:derivedFrom}).
  \item No forced ALIVE (both ALIVE and PARTIAL verdict paths documented).
  \item Zero \texttt{.ggen} source files in \texttt{prompt-manufactory/ggen/}.
  \item All 22 legacy \texttt{.ggen} files classified with remediation routes.
\end{enumerate}

Two gates remain pending: PI\_INTEL topology completion (dependent on the concurrent
\texttt{PI\_RESEARCH\_PROGRAM\_INTEL\_001} workflow completing its classification phase),
and implementation of the remaining 7 Tera templates. The current verdict is correctly
PARTIAL.
